\documentclass[9pt]{article}
% Kiln CV template: a white drawing sheet, IBM Plex, ink rules, one ember thing (the mark).
% The alternative Kiln-style CV (kept for later; the main CV is cv/CV_Borodin_EN.tex). Shared by cv/kiln/CV_Borodin_EN.tex and CV_Borodin_RU.tex. Build from the repo root with
%   xelatex -output-directory=cv/build cv/kiln/CV_Borodin_EN.tex
% Fonts: IBM Plex Sans Condensed and IBM Plex Mono from IBM's own release (Cyrillic included).

\usepackage[a4paper,top=10mm,bottom=7mm,left=12mm,right=12mm]{geometry}
\usepackage{fontspec}
\usepackage{xcolor}
\usepackage{graphicx}
\usepackage{tikz}
\usepackage{enumitem}
\usepackage{tabularx}
\usepackage{array}
\usepackage[hidelinks,unicode]{hyperref}

% ---------------------------------------------------------------- tokens (Kiln 4 · Colour)
\definecolor{ink}{HTML}{16130F}
\definecolor{ink2}{HTML}{605A52}
\definecolor{dimmed}{HTML}{9F998F}
\definecolor{rule}{HTML}{D1CDC6}
\definecolor{ember}{HTML}{E84C14}

% ---------------------------------------------------------------- type (Kiln 3 · Type)
\setmainfont{IBM Plex Sans Cond}[BoldFont={IBM Plex Sans Cond SmBld}, ItalicFont={IBM Plex Sans Cond}]
\newfontfamily\monofont{IBM Plex Mono Medm}[BoldFont={IBM Plex Mono SmBld}]
\newfontfamily\labelfont{IBM Plex Mono Medm}[LetterSpace=11]

\color{ink}
\pagestyle{empty}
\setlength{\parindent}{0pt}
\setlength{\parskip}{0pt}
\setlength{\tabcolsep}{0pt}
\emergencystretch=2em
\raggedright

\hyphenpenalty=10000 \exhyphenpenalty=10000   % ragged right, no broken words (Kiln 3 · Type)

% the label column, used where a label names a row (skills)
\newlength{\labw}\setlength{\labw}{25mm}
\newlength{\bodyw}\setlength{\bodyw}{\dimexpr\textwidth-\labw\relax}

% ---------------------------------------------------------------- pieces
% Mono capitals with +0.11em, the label voice
\newcommand{\lab}[1]{{\labelfont\fontsize{6.6}{9}\selectfont\color{ink2}\MakeUppercase{#1}}}
% Dates and small technical facts in mono
\newcommand{\mono}[1]{{\monofont\fontsize{7.4}{10}\selectfont\color{ink2}#1}}
% A link shown as its address, quiet
\newcommand{\addr}[2]{\href{#1}{\mono{#2}}}

% Section: a heavy ink rule, then the label (the site's section pattern)
\newcommand{\cvsection}[1]{%
  \par\vspace{3pt}%
  {\color{ink}\rule{\textwidth}{1.4pt}}\par
  \vspace{2.5pt}\lab{#1}\par\vspace{3pt}%
}

% An entry: title left, date right in mono (the site's job row), place under it, then the list
\newcommand{\entry}[4]{% #1 date, #2 title, #3 place, #4 list
  \par\noindent{\bfseries\fontsize{9.8}{12}\selectfont #2}\hfill\mono{#1}\par
  {\color{ink2}#3}\par
  #4\par\vspace{3pt}%
}

% A labelled row: mono label left, text right (\leavevmode keeps both first baselines aligned)
\newcommand{\row}[2]{%
  \par\noindent
  \parbox[t]{\labw}{\leavevmode\raggedright\mono{#1}}%
  \parbox[t]{\dimexpr\linewidth-\labw\relax}{\leavevmode\raggedright #2}\par\vspace{2.5pt}%
}

% A one-line fact with its date on the right
\newcommand{\fact}[2]{\par\noindent #1\hfill\mono{#2}\par\vspace{1pt}}

\setlist[itemize]{leftmargin=1em,labelsep=0.35em,itemsep=0pt,topsep=1.5pt,parsep=0pt,partopsep=0pt}
\renewcommand\labelitemi{{\color{dimmed}—}}

% The photo, cropped square with Kiln's panel radius
\newcommand{\photo}[2]{%
  \begin{tikzpicture}
    \clip[rounded corners=2.2mm] (0,0) rectangle (#2,#2);
    \node[anchor=south west,inner sep=0] at (0,0) {\includegraphics[width=#2,height=#2]{#1}};
  \end{tikzpicture}%
}

% Header: lockup left, contacts right, the heavy rule under it (the site header)
\newcommand{\cvheader}[2]{% #1 lockup pdf, #2 contact lines
  \noindent
  \begin{minipage}[c]{0.5\textwidth}\includegraphics[height=11mm]{#1}\end{minipage}%
  \begin{minipage}[c]{0.5\textwidth}\raggedleft #2\end{minipage}\par
  \vspace{5pt}{\color{ink}\rule{\textwidth}{1.4pt}}\par\vspace{7pt}%
}

% Lede size (Russian sets wider, so its file may override this)
\newcommand{\ledesize}{\fontsize{10}{12.6}\selectfont}

% Title block: label, name, lede; photo on the right
\newcommand{\cvtitle}[4]{% #1 label, #2 name, #3 lede, #4 photo file
  \noindent
  \begin{minipage}[t]{\dimexpr\textwidth-27mm\relax}
    \raggedright\lab{#1}\par\vspace{3pt}
    {\bfseries\fontsize{22}{25}\selectfont #2}\par\vspace{3pt}
    {\ledesize\color{ink2}#3}
  \end{minipage}\hfill
  \begin{minipage}[t]{22mm}\vspace{0pt}\photo{#4}{22mm}\end{minipage}\par
}

\fontsize{9}{11}\selectfont

% Two columns of labelled rows (short values, e.g. skills)
\newenvironment{twocol}{\par\noindent\begin{minipage}[t]{0.485\textwidth}}{\end{minipage}\par}
\newcommand{\nextcol}{\end{minipage}\hfill\begin{minipage}[t]{0.485\textwidth}}

\hypersetup{
  pdftitle={Alexander Borodin -- ML Engineer | LLMs \& AI Agents},
  pdfauthor={Alexander Borodin},
  pdfsubject={Curriculum vitae}
}

\begin{document}
\fontsize{9}{11}\selectfont

\cvheader{cv/kiln/lockup.pdf}{%
  \href{mailto:apdorodin765@gmail.com}{\mono{apdorodin765@gmail.com}}\quad
  \mono{+7 908 971 16 97}\par
  \addr{https://github.com/aleksanderborodin}{github.com/aleksanderborodin}\quad
  \addr{https://t.me/aleksanderbor}{t.me/aleksanderbor}\par
  \addr{https://aleksanderborodin.github.io/}{aleksanderborodin.github.io}}

\cvtitle{ML Engineer · LLMs and AI agents}{Alexander Borodin}%
{ML engineer and Applied Mathematics and Physics student at MIPT. I build LLM systems, from RAG and model integrations to multi-agent workflows and agent-driven evolution: iterative search and improvement of solutions. My experience includes benchmark development, model evaluation, and applied optimization.}%
{assets/profile-original.jpg}

\cvsection{Experience}
\entry{Feb–Aug 2026 · 6 months}{ML Engineer — Research Internship}{CKM, MIPT}{%
\begin{itemize}
  \item Developed LLM-based agents and benchmarks to evaluate their performance.
  \item Worked on agent-driven evolution: systems in which agents iteratively generate, evaluate, and improve solutions.
  \item Contributed to two research papers: one under review at ACL and another in preparation for submission.
\end{itemize}}
\entry{Jan–May 2026}{Product / AI Engineer}{\href{https://modelgate.ru}{ModelGate.ru} — LLM API gateway and proxy service}{%
\begin{itemize}
  \item Developed an LLM API gateway providing a unified interface to models from multiple providers.
  \item Integrated provider APIs and implemented request routing with model and provider selection.
  \item Built core service components, including the backend and public website. Stack: Python, LLM APIs.
\end{itemize}}
\entry{Sep 2024 – Feb 2025}{Model and Algorithm Developer}{Regional Bank, client project}{%
\begin{itemize}
  \item Developed a cash collection optimization algorithm; historical-data analysis projected a 20M RUB annual profit increase.
  \item Built a CBR web-scraping/analysis pipeline for data unavailable via API and an internal document Q\&A RAG system.
  \item Stack: Python, NumPy, Pandas, RAG, web scraping.
\end{itemize}}

\cvsection{Projects}
\entry{2026}{Idea Evolve}{\addr{https://github.com/aleksanderborodin/idea_evolve}{github.com/aleksanderborodin/idea\_evolve}}{%
\begin{itemize}
  \item Evolutionary multi-agent framework: specialized Claude agents work in parallel across generations and build a shared knowledge base.
  \item Improved Sidon set search from 66 to 105 elements in 7 autonomous generations, passing the √N ≈ 100 reference for N = 10,000.
  \item Implemented role-specialized agents, resumable filesystem state, evaluation/knowledge extraction, and a Flask dashboard.
  \item Stack: Python 3.12+, Claude Code CLI, OpenCode, Anthropic API, Flask, multi-agent orchestration.
\end{itemize}}
\entry{2025}{Prefix Injection Jailbreak}{\addr{https://github.com/aleksanderborodin/prefix_injection_jailbreak}{github.com/aleksanderborodin/prefix\_injection\_jailbreak}}{%
\begin{itemize}
  \item Researched Choice-Blindness-inspired prefix attacks; templates made Gemini 2.5 Pro and DeepSeek R1 complete MaliciousInstruct prompts with near-100\% success vs. a near-zero baseline.
  \item Stack: Python, LLM APIs, model evaluation, prompt/context engineering.
\end{itemize}}
\entry{2025}{Yoola Explain}{\addr{https://github.com/aleksanderborodin/yoola_explain}{github.com/aleksanderborodin/yoola\_explain}}{%
\begin{itemize}
  \item Browser extension that extracts Terms of Service and returns structured LLM summaries: key points, data collection notes, and warnings.
  \item Added multi-language support and caching. Stack: JavaScript, FastAPI, OpenRouter, LLM APIs.
\end{itemize}}

\cvsection{Education}
\entry{2023–2027}{Bachelor's degree, Applied Mathematics and Physics}{FPMI MIPT, in progress · GPA 7.8/10}{%
Coursework: mathematical analysis, probability, linear algebra, algorithms, computational complexity, ML, physics, economics.}

\cvsection{Skills}
\begin{twocol}
\row{Programming}{Python advanced; C++ intermediate; C and Assembler basic code reading.}
\row{Data / ML}{PyTorch basic; NumPy; Pandas; PostgreSQL; Matplotlib; ClearML; Excel.}
\row{NLP / LLM}{BM25; embeddings; RAG; LLM APIs; evaluation; prompt/context engineering; jailbreak testing.}
\nextcol
\row{Agentic tools}{Claude Code CLI; OpenCode; multi-agent orchestration; filesystem-based workflows.}
\row{Backend}{FastAPI; Flask; LangChain; LangGraph; Linux home server administration and workstation usage.}
\row{Languages}{Russian native; English B2+ certified by MIPT placement test.}
\end{twocol}

\cvsection{Programs and achievements}
\entry{Jul 21 – Aug 4, 2026}{\href{https://airi.net/ru/summer-school-2026/}{Summer with AIRI 2026} — Artificial Intelligence Summer School}{Final team project: \href{https://github.com/n4rly-boop/ideas-lake}{Ideas Lake}, long-term memory for evolutionary AI workflows using hybrid retrieval and a knowledge graph.}{}
\fact{\textbf{Gold Medal}, \href{https://iepho.ru/}{International Experimental Physics Olympiad (IEPhO)}}{2023}
\fact{\textbf{Prize-winner}, \href{https://vos.olimpiada.ru/}{All-Russian School Olympiad in Physics}, final stage — Grades 10 and 11}{2022, 2023}

\end{document}
